\subsection{小问1：模型建立、求解与验证}

\subsubsection{任务目标与模型口径}
说明本小问要求的输出、单位、粒度、验收标准及求解属性。

\subsubsection{数据、假设与符号}
说明输入字段、数据质量、关键假设以及符号与量纲。

\subsubsection{从题意到数学模型}
依次给出基础模型、约束或边界条件、参数来源与完整模型。

\subsubsection{求解方法与实现}
说明基线、最终算法、核心/辅助分类及理由、随机种子和求解设置。若本小问包含核心算法，在完整公式和参数来源后使用 algorithm2e 给出带行号、输入、输出、循环/分支、终止与 fallback 的伪代码，再插入图形流程图，并解释二者与真实代码和输出字段的对应关系。流程图源文件保存在本小问“代码”目录，渲染文件保存在“图”目录。

\subsubsection{结果与验证}
\label{sec:q1-s1-validation}
引用本小问结果目录中的核心结果表。对每一种验证评估方法，依次说明验证对象、选择理由、必要定义、本题实际协议、指标或判据、真实结果及定位、结果解释和决策影响；不得写成与当前模型和结果脱节的通用介绍。

\subsubsection{直接回答}
用与结果文件一致的口径直接回答本小问，并说明适用边界。
